%
% Copyright (C) 1995, 2000, 2013 American Mathematical Society.
% Copyright (C) 2016-2025 LaTeX Project and American Mathematical Society.
%
% This work may be distributed and/or modified under the
% conditions of the LaTeX Project Public License, either version 1.3c
% of this license or (at your option) any later version.
% The latest version of this license is in
%   https://www.latex-project.org/lppl.txt
% and version 1.3c or later is part of all distributions of LaTeX
% version 2005/12/01 or later.
%
% This work has the LPPL maintenance status `maintained'.
%
% The Current Maintainer of this work is the LaTeX Project.
%
\documentclass[leqno,titlepage,openany]{amsldoc}[1999/12/13]

\def\MaintainedByLaTeXTeam#1{%
\begin{center}%
\fbox{\fbox{\begin{tabular}{@{}l@{}}%
This file is maintained by the \LaTeX{} Project team.\\%
Bug reports can be opened (category \texttt{#1}) at\\%
\url{https://latex-project.org/bugs/}.\end{tabular}}}\end{center}}

\ifx\UndEfiNed\url
  \ClassError{amsldoc}{%
    This version of amsldoc.tex must be processed\MessageBreak
    with a newer version of amsldoc.cls (2.02 or later).}{}
\fi

\renewcommand\bibname{References}

%    Definitions that should be moved to amsldoc.cls
%    when it is next updated
\makeatletter
\let\cleardouble@page\cleardoublepage
\AtBeginDocument{%
  \ifx\cleardouble@page\cleardoublepage
    \def\cleardoublepage{\clearpage{\pagestyle{empty}\cleardouble@page}}
  \fi
}

% #1 is empty or one of class, environment, option, package
% could use different font in each case
\makeatletter
\def\@category@index#1#2{%
    \autoindex{%
        #2\actualchar\string\index#1font{#2}%
        \ifx\@nil#1\@nil\else\space#1\fi
    }%
}

\let\indexpackagefont\textsf
\let\indexclassfont\textsf
\let\indexoptionfont\texttt
\let\indexenvironmentfont\texttt
\let\indexfont\texttt
\makeatother
%    For a package that won't be indexed, i.e. amsmath
\newcommand{\nipkg}{\textsf}
\def\tikz/{Ti\textit{k}Z}
%    End definitions for amsldoc.cls

\title{User's Guide for the \nipkg{amsmath} Package (Version~2.1)}
\author{American Mathematical Society, \LaTeX\ Project}
\date{1999-12-13\\(revised 2002-02-25, 2016-11-14, 2018-04-05, 2019-10-14, 2020-02-18)}
\makeatletter
\def\@thanks{\bigskip\MaintainedByLaTeXTeam{amslatex}}
\makeatother

% Very minor changes are not listed in the \date above:
%
% 2021-08-22: reported typo fixed mathools -> mathtools
% 2025-05-23: drop spurious option list


%    Use the amsmath package and amscd package in order to print
%    examples.
\usepackage{amsmath}
\usepackage{amscd}

\usepackage{imakeidx}
%    Need alternate coding for vertical bar to get it into index.
%    Based on \qcbang and \cnbang in amsldoc.cls, refined by David Carlisle.
%    Give it a sort field that will be at the end of the symbols;
%    must avoid " and | which have special meanings.
\DeclareRobustCommand{\qcvert}{\qc\|\index{{z1}@{\ntt\char'174}}}
\DeclareRobustCommand{\cnvert}{%
  \ncn{\|}\index{{z0}@{\ntt\bslchar\char'174}}}
%    Need alternate index sort field for bang to allow & to sort first
\DeclareRobustCommand{\altcnbang}{%
  \ncn{\!}\index{"!b@{\ntt\bslchar\qcbang}}}
\DeclareRobustCommand{\cnamp}{%
  \index{"!a@\texttt{\&}}}

\makeindex[intoc] % generate index data
\providecommand{\see}[2]{\textit{see} #1}

%    The amsldoc class includes a number of features useful for
%    documentation about TeX, including:
%
%    ---Commands \tex/, \amstex/, \latex/, ... for uniform treatment
%    of the various logos and easy handling of following spaces.
%
%    ---Commands for printing various common elements: \cn for command
%    names, \fn for file names (including font-file names), \env for
%    environments, \pkg and \cls for packages and classes, etc.

%    Many of the command names used here are rather long and will
%    contribute to poor linebreaking if we follow the \latex/ practice
%    of not hyphenating anything set in tt font; instead we selectively
%    allow some hyphenation.
\allowtthyphens % defined in amsldoc.cls

\hyphenation{ac-cent-ed-sym-bol add-to-counter add-to-length align-at
  aligned-at allow-dis-play-breaks ams-art ams-cd ams-la-tex amsl-doc
  ams-symb ams-tex ams-text ams-xtra bmatrix bold-sym-bol cen-ter-tags
  eqn-ar-ray idots-int int-lim-its latex med-space neg-med-space
  neg-thick-space neg-thin-space no-int-lim-its no-name-lim-its
  over-left-arrow over-left-right-arrow over-right-arrow pmatrix
  qed-sym-bol set-length side-set small-er tbinom the-equa-tion
  thick-space thin-space un-der-left-arrow un-der-left-right-arrow
  un-der-right-arrow use-pack-age var-inj-lim var-proj-lim vmatrix
  xalign-at xx-align-at}

%    Prepare for illustrating the \vec example
\newcommand{\vect}[1]{\mathbf{#1}}

\newcommand{\booktitle}[1]{\textit{#1}}
\newcommand{\journalname}[1]{\textit{#1}}
\newcommand{\seriesname}[1]{\textit{#1}}

%    Command to insert and index a particular phrase. Doesn't work for
%    certain kinds of special characters in the argument.
\newcommand{\ii}[1]{#1\index{#1}}

\newcommand{\vstrut}[1]{\vrule width0pt height#1\relax}

%    An environment for presenting comprehensive address information:
\newenvironment{infoaddress}{%
  \par\topsep\medskipamount
  \trivlist\centering
  \item[]%
  \begin{minipage}{.7\columnwidth}%
  \raggedright
}{%
  \end{minipage}%
  \endtrivlist
}

\newenvironment{eqxample}{%
  \par\addvspace\medskipamount
  \noindent\begin{minipage}{.5\columnwidth}%
  \def\producing{\end{minipage}\begin{minipage}{.5\columnwidth}%
    \hbox\bgroup\kern-.2pt\vrule width.2pt%
      \vbox\bgroup\parindent0pt\relax
%    The 3pt is to cancel the -\lineskip from \displ@y
    \abovedisplayskip3pt \abovedisplayshortskip\abovedisplayskip
    \belowdisplayskip0pt \belowdisplayshortskip\belowdisplayskip
    \noindent}
}{%
  \par
%    Ensure that a lonely \[\] structure doesn't take up width less than
%    \hsize.
  \hrule height0pt width\hsize
  \egroup\vrule width.2pt\kern-.2pt\egroup
  \end{minipage}%
  \par\addvspace\medskipamount
}

%    Errors in output shouldn't be \texttt, and they may also be longer
%    than one line, but they should be identified by the "errorbullet".
%    (They were formerly labeled as aubsections.)
\newenvironment{erroro}[1]{%
  \par\addvspace\medskipamount
  \noindent\hangindent\parindent
  \hbox to\parindent{\errorbullet\hfil}\ignorespaces
  #1\par\vspace{\smallskipamount}
}{%
  \par\addvspace\medskipamount
}

%    The chapters are so short, perhaps we shouldn't call them by the
%    name `Chapter'. We make \chaptername read an argument in order to
%    remove a following \space or "{} " (both possibilities are present
%    in book.cls).

\renewcommand{\chaptername}[1]{}
\newcommand{\chapnum}[1]{\mdash #1\mdash }
\makeatletter
\def\@makechapterhead#1{%
  \vspace{1.5\baselineskip}%
  {\parindent \z@ \raggedright \reset@font
    \ifnum \c@secnumdepth >\m@ne
      \large\bfseries \chapnum{\thechapter}%
      \par\nobreak
      \vskip.5\baselineskip\relax
    \fi
    #1\par\nobreak
    \vskip\baselineskip
  }}
\makeatother

%    A command for ragged-right parbox in a tabular.
\newcommand{\rp}{\let\PBS\\\raggedright\let\\\PBS}

%    Non-indexed file name
\newcommand{\nfn}[1]{\texttt{#1}}

%    For the examples in the math spacing table.
%%\newcommand{\lspx}{\mbox{\rule{5pt}{.6pt}\rule{.6pt}{6pt}}}
%%\newcommand{\rspx}{\mbox{\rule[-1pt]{.6pt}{7pt}%
%%  \rule[-1pt]{5pt}{.6pt}}}
\newcommand{\lspx}{\mathord{\Rightarrow\mkern-1mu}}
\newcommand{\rspx}{\mathord{\mkern-1mu\Leftarrow}}
\newcommand{\spx}[1]{$\lspx #1\rspx$}

%    For a list of characters representing document input.
\newcommand{\clist}[1]{%
  \mbox{\ntt\spaceskip.2em plus.1em \xspaceskip\spaceskip#1}}

%    Fix weird \latex/ definition of rightmark.
\makeatletter
\def\rightmark{\expandafter\@rightmark\botmark{}{}}
%    Also turn off section marks.
\let\sectionmark\@gobble
\renewcommand{\chaptermark}[1]{%
  \uppercase{\markboth{\rhcn#1}{\rhcn#1}}}
\newcommand{\rhcn}{\thechapter. }
\makeatother

%    Include down to \section but not \subsection, in toc:
\setcounter{tocdepth}{1}

%    Subheadings for the bibliography
\newif\iffirstbibsubhead \firstbibsubheadtrue
\newcommand{\bibsubhead}[1]{%
  \iffirstbibsubhead \firstbibsubheadfalse
  \else \addvspace{\medskipamount}
  \fi
  \item[]\hspace*{-\leftmargin}\textbf{#1}\par
  \vspace{\smallskipamount}
}

\DeclareMathOperator{\ix}{ix}
\DeclareMathOperator{\nul}{nul}
\DeclareMathOperator{\End}{End}
\DeclareMathOperator{\xxx}{xxx}

\usepackage{hyperref}
\hypersetup{pdfborder=0 0 0}
\makeatletter
\let\oldcs\cs
\def\cs#1{\texorpdfstring{\oldcs{#1}}{\@backslashchar\@backslashchar#1}}
\let\cn\cs
\makeatother
\begin{document}
\section{Error messages}

\begin{error}{\begin{split} won't work here.}
\errexa
\begin{verbatim}
! Package amsmath Error: \begin{split} won't work here.
 ...

l.8 \begin{split}

? h
\Did you forget a preceding \begin{equation}?
If not, perhaps the `aligned' environment is what you want.
?
\end{verbatim}
\errexpl
The \env{split} environment does not construct a stand-alone displayed
equation; it needs to be used within some other environment such as
\env{equation} or \env{gather}.

\end{error}

\begin{error}{Erroneous nesting of equation structures}
\errexa
\begin{verbatim}
! Package amsmath Error: Erroneous nesting of equation structures;
(amsmath)                trying to recover with `aligned'.

See the amsmath package documentation for explanation.
Type  H <return>  for immediate help.
 ...

l.260 \end{alignat*}
                    \end{equation*}
\end{verbatim}
\errexpl
The structures \env{align}, \env{alignat}, etc., are designed
for top-level use and for the most part cannot be nested inside some
other displayed equation structure. The chief exception is that
\env{align} and most of its variants can be used inside the
\env{gather} environment.
\end{error}

\begin{error}{Extra & on this line}
\errexa
\begin{verbatim}
! Package amsmath Error: Extra & on this line.

See the amsmath package documentation for explanation.
Type  H <return>  for immediate help.
 ...

l.9 \end{alignat}

? h
\An extra & here is so disastrous that you should probably exit
 and fix things up.
?
\end{verbatim}
\errexpl
In an \env{alignat} structure the number of alignment points per line
is dictated by the numeric argument given after |\begin{alignat}|.
If you use more alignment points in a line it is assumed that you
accidentally left out a newline command \cn{\\} and the above error is
issued.
\end{error}

%% **********
\enlargethispage{1.5\baselineskip}

\begin{error}{Font OMX/cmex/m/n/7=cmex7 not loadable ...}
\errexa
\begin{verbatim}
! Font OMX/cmex/m/n/7=cmex7 not loadable: Metric (TFM) file not found.
<to be read again>
                   relax
l.8 $a
      b+b^2$
? h
I wasn't able to read the size data for this font,
so I will ignore the font specification.
[Wizards can fix TFM files using TFtoPL/PLtoTF.]
You might try inserting a different font spec;
e.g., type `I\font<same font id>=<substitute font name>'.
?
\end{verbatim}
\errexpl
Certain extra sizes of some Computer Modern fonts that were formerly
available mainly through the AMSFonts\index{amsfont@AMSFonts collection}
distribution are considered part of standard \latex/ (as of June 1994):
\fn{cmex7}\ndash \texttt{9}, \fn{cmmib5}\ndash \texttt{9}, and
\fn{cmbsy5}\ndash \texttt{9}. If these extra sizes are missing on your
system, you should try first to get them from the source where you
obtained \latex/. If that fails, you could try getting the fonts from
CTAN (e.g., in the form of Metafont\index{Metafont source files} source
files, directory \nfn{/tex-archive/fonts/latex/mf}, or in PostScript
Type 1 format, directory
\nfn{/tex-archive/fonts/cm/ps-type1/bakoma}\index{BaKoMa fonts}\relax
\index{PostScript fonts}).

If the font name begins with \fn{cmex}, there is a special option%
\index{options!behavior of particular options}
\opt{cmex10} for the \nipkg{amsmath} package that provides a temporary
workaround. I.e., change the \cn{usepackage} to
\begin{verbatim}
\usepackage[cmex10]{amsmath}
\end{verbatim}
This will force the use of the 10-point size of the \fn{cmex} font in
all cases. Depending on the contents of your document this may be
adequate.
\end{error}

\enlargethispage{1\baselineskip}
\begin{error}{Improper argument for math accent}
\errexa
\begin{verbatim}
! Package amsmath Error: Improper argument for math accent:
(amsmath)                Extra braces must be added to
(amsmath)                prevent wrong output.

See the amsmath package documentation for explanation.
Type  H <return>  for immediate help.
 ...

l.415 \tilde k_{\lambda_j} = P_{\tilde \mathcal
                                               {M}}
?
\end{verbatim}
\errexpl
Non-simple arguments for any \LaTeX{} command should be enclosed in
braces. In this example extra braces are needed as follows:
\begin{verbatim}
... P_{\tilde{\mathcal{M}}}
\end{verbatim}
\end{error}

\begin{error}{Math formula deleted: Insufficient extension fonts}
\errexa
\begin{verbatim}
! Math formula deleted: Insufficient extension fonts.
l.8 $ab+b^2$

?
\end{verbatim}
\errexpl
This usually follows a previous error |Font ... not loadable|; see the
discussion of that error (above) for solutions.
\end{error}

\enlargethispage{1\baselineskip}
\begin{error}{Missing number, treated as zero}
\errexa
\begin{verbatim}
! Missing number, treated as zero.
<to be read again>
                   a
l.100 \end{alignat}

? h
A number should have been here; I inserted `0'.
(If you can't figure out why I needed to see a number,
look up `weird error' in the index to The TeXbook.)

?
\end{verbatim}
\errexpl
There are many possibilities that can lead to this error. However, one
possibility that is relevant for the \nipkg{amsmath} package is that you
forgot to give the number argument of an \env{alignat} environment, as
in:
\begin{verbatim}
\begin{alignat}
 a&  =b&    c& =d\\
a'& =b'&   c'& =d'
\end{alignat}
\end{verbatim}
where the first line should read instead
\begin{verbatim}
\begin{alignat}{2}
\end{verbatim}

Another possibility is that you have a left bracket character |[|
following a line break\index{line break} command \cn{\\} in a multiline
construction such as \env{array}, \env{tabular}, or \env{eqnarray}.
This will be interpreted by \latex/ as the beginning of an `additional
vertical space request \cite[\S C.1.6]{lamport}, even if it occurs on the
following line and is intended to be part of the contents. For example,
\begin{verbatim}
\begin{array}{c}
a+b\\
[f,g]\\
m+n
\end{array}
\end{verbatim}
To prevent the error message in such a case, you can
add braces as discussed in the \latex/ manual \cite[\S C.1.1]{lamport}:
\begin{verbatim}
\begin{array}{c}
a+b\\
{[f,g]}\\
m+n
\end{array}
\end{verbatim}
or precede the bracketed expression by \cn{relax}.
\end{error}

\goodbreak

\begin{error}{Missing \right. inserted}
\errexa
\begin{verbatim}
! Missing \right. inserted.
<inserted text>
                \right .
l.10 \end{multline}

? h
I've inserted something that you may have forgotten.
(See the <inserted text> above.)
With luck, this will get me unwedged. But if you
really didn't forget anything, try typing `2' now; then
my insertion and my current dilemma will both disappear.
\end{verbatim}
\errexpl
This error typically arises when you try to insert a
line break\index{line break} or |&|\cnamp\ between a
\cn{left}-\cn{right} pair of delimiters\index{delimiters} in any
multiline environment, including \env{split}.
\begin{verbatim}
\begin{multline}
AAA\left(BBB\\
  CCC\right)
\end{multline}
\end{verbatim}
There are two possible solutions: (1)~instead of using \cn{left} and
\cn{right}, use `big' delimiters\index{delimiters!fixed size} of fixed
size (\cn{bigl} \cn{bigr} \cn{biggl} \cn{biggr} \dots; see
\secref{bigdel}); or (2)~use null delimiters\index{delimiters!null} to
break up the \cn{left}-\cn{right} pair into parts for each line (or cell):
\begin{verbatim}
AAA\left(BBB\right.\\
  \left.CCC\right)
\end{verbatim}
The latter solution may result in mismatched
delimiter\index{delimiters!mismatched sizes} sizes;
ensuring that they match requires using \cn{vphantom} in the segment
that has the smaller delimiter (or possibly \cn{smash} in the segment
that has the larger delimiter). In the argument of \cn{vphantom} put a
copy of the tallest element that occurs in the other segment, e.g.,
\begin{verbatim}
xxx \left(\int_t yyy\right.\\
  \left.\vphantom{\int_t} zzz ... \right)
\end{verbatim}
\end{error}

\goodbreak

\begingroup
\lccode`M=`M
\lccode`@=`\}
\lowercase{\endgroup
\begin{error}{Missing @ inserted}}

\errexa
\begin{verbatim}
! Missing } inserted.
<inserted text>
                \right .
l.10 \end{multline}

? h
I've inserted something that you may have forgotten.
(See the <inserted text> above.)
With luck, this will get me unwedged. But if you
really didn't forget anything, try typing `2' now; then
my insertion and my current dilemma will both disappear.
\end{verbatim}
\errexpl
This error is often the result of using \verb+$+ within a multiline
display environment.  Remove such \verb+$+ signs (except when they
appear within \verb+\text{...}+).
\end{error}

\begin{error}{Old form `\pmatrix' should be \begin{pmatrix}.}
\errexa
\begin{verbatim}
! Package amsmath Error: Old form `\pmatrix' should be
                         \begin{pmatrix}.

See the amsmath package documentation for explanation.
Type  H <return>  for immediate help.
 ...

\pmatrix ->\left (\matrix@check \pmatrix
                                         \env@matrix
l.16 \pmatrix
             {a&b\cr c&d\cr}
? h
`\pmatrix{...}' is old Plain-TeX syntax whose use is
ill-advised in LaTeX.
?
\end{verbatim}
\errexpl
When the \nipkg{amsmath} package is used, the old forms of \cn{pmatrix},
\cn{matrix}, and \cn{cases} cannot be used any longer because of naming
conflicts. Their syntax did not conform with standard \LaTeX{} syntax
in any case.
\end{error}

\enlargethispage{2\baselineskip}
\begin{error}{Paragraph ended before \xxx was complete}
\errexa
\begin{verbatim}
Runaway argument?

! Paragraph ended before \multline was complete.
<to be read again>
                   \par
l.100

? h
I suspect you've forgotten a `}', causing me to apply this
control sequence to too much text. How can we recover?
My plan is to forget the whole thing and hope for the best.
?
\end{verbatim}
\errexpl
This might be produced by a blank line between the |\begin| and |\end|.
Another possibility is a misspelling in the |\end{multline}| command,
e.g.,
\begin{verbatim}
\begin{multline}
...
\end{multiline}
\end{verbatim}
or by using abbreviations for certain environments, such as |\bal| and
|\eal| for |\begin{align}| and |\end{align}|:
\begin{verbatim}
\bal
...
\eal
\end{verbatim}
For technical reasons that kind of abbreviation does not work with
the more complex displayed equation environments of the \nipkg{amsmath} package
(\env{gather}, \env{align}, \env{split}, etc.; cf.\ \fn{technote.tex}).
\end{error}

\begin{error}{Runaway argument?}
See the discussion for the error message
\texttt{Paragraph ended before \ncn{xxx} was complete}.
\end{error}

\begin{error}{Unknown option `xxx' for package `yyy'}
\errexa
\begin{verbatim}
! LaTeX Error: Unknown option `intlim' for package `amsmath'.
...
? h
The option `intlim' was not declared in package `amsmath', perhaps you
misspelled its name. Try typing  <return>  to proceed.
?
\end{verbatim}
\errexpl
This means that you misspelled the option\index{options!unknown}
name, or the package simply does not have an option that you expected
it to have. Consult the documentation for the given package.
\end{error}

%% ------------------------------------------------------------------ %%

\section{Warning messages}

\begin{error}{Cannot use `split' here}
\errexa
\begin{verbatim}
Package amsmath Warning: Cannot use `split' here;
(amsmath)                trying to recover with `aligned'
\end{verbatim}
\errexpl The \env{split} environment is designed to serve as the entire
body of an equation, or an entire line of an \env{align} or \env{gather}
environment. There cannot be any printed material before or
after it within the same enclosing structure:
\begin{verbatim}
\begin{equation}
\left\{ % <-- Not allowed
\begin{split}
...
\end{split}
\right. % <-- Not allowed
\end{equation}
\end{verbatim}
\end{error}

\begin{error}{Foreign command \over [or \atop or \above]}
\errexa
\begin{verbatim}
Package amsmath Warning: Foreign command \over; \frac or \genfrac
(amsmath)                should be used instead.
\end{verbatim}
\errexpl The primitive generalized fraction commands of \tex/\mdash
\cs{over}, \cs{atop}, \cs{above}\mdash are deprecated when the
\nipkg{amsmath} package is used because their syntax is foreign to \latex/
and \nipkg{amsmath} provides native \latex/ equivalents. See
\fn{technote.tex} for further information.
\end{error}

%% ------------------------------------------------------------------ %%

\section{Wrong output}

\begin{erroro}{Section numbers 0.1, 5.1, 8.1 instead of 1, 2, 3}
%
This most likely means that you have the arguments for \cn{numberwithin}
in reverse order:
\begin{verbatim}
\numberwithin{section}{equation}
\end{verbatim}
That means `print the section number as \textit{equation
number}.\textit{section number} and reset to 1 every time an equation
occurs' when what you probably wanted was the inverse
\begin{verbatim}
\numberwithin{equation}{section}
\end{verbatim}
\end{erroro}

\begin{erroro}{The \cn{numberwithin} command had no effect on equation
  numbers}
%
Are you looking at the first section in your document? Check the section
numbers elsewhere to see if the problem is the one described in the
previous item.
\end{erroro}

\begin{erroro}{Bracketed material disappears at the beginning of
  \env{aligned} or \env{gathered}}
%
For most multiline equation environments, \nipkg{amsmath} disables the
\latex/ convention to recognize a [bracketed] expression as an option%
\index{ed env@\texttt{-ed} environments}%
\index{options!positioning of \texttt{-ed} environments}%
\index{options!space before \texttt{[}}
when it occurs at the beginning of an environment or immediately
following \cn{\\}, provided a space intervenes.  This fails at the
beginning of \env{aligned} and \env{gathered}, and is a bug.  Insert
\cn{relax} before the opening brace to restore the desired result.
\end{erroro}

%%%%%%%%%%%%%%%%%%%%%%%%%%%%%%%%%%%%%%%%%%%%%%%%%%%%%%%%%%%%%%%%%%%%%%%%

\end{document}
